\documentclass[a4paper,11pt]{article}
\usepackage[utf8]{inputenc}
\usepackage[spanish]{babel}
\usepackage{geometry}
\geometry{a4paper, margin=2.5cm}
\usepackage{graphicx}
\usepackage{hyperref}
\usepackage{booktabs}
\usepackage{xcolor}
\usepackage{colortbl}
\usepackage{tikz}
\usetikzlibrary{shapes.geometric, arrows, positioning, babel}
\usepackage{amsmath}

\title{\textbf{Reporte Científico de Modelado de Datos: Análisis Exhaustivo} \\ \Large BCP Datafest 2026}
\author{Aron}
\date{\today}

\begin{document}

\maketitle
\tableofcontents
\newpage

\section{Resumen Ejecutivo}
El presente informe documenta el proceso iterativo y la resolución científica del reto de modelado predictivo para el BCP Datafest. El objetivo fue maximizar la métrica AUC identificando clientes propensos a la compra ("1"), superando la barrera del desbalanceo poblacional (85\% "0" vs 15\% "1") y el severo solapamiento de clases descubierto mediante análisis espacial.

\section{Evolución de Modelos (Base vs. "Locuras")}

La investigación se dividió en dos fases: los modelos clásicos (urón, baseline) y los enfoques no convencionales o "locuras". A continuación, se presenta la matriz comparativa exhaustiva, desglosando el rendimiento en la Clase 0 (Mayoritaria) y Clase 1 (Minoritaria).

\begin{table}[h!]
\centering
\resizebox{\textwidth}{!}{
\begin{tabular}{l c ccc ccc c}
\toprule
 & & \multicolumn{3}{c}{\textbf{Clase 0 (No Compra)}} & \multicolumn{3}{c}{\textbf{Clase 1 (Compra)}} & \\
\cmidrule(lr){3-5} \cmidrule(lr){6-8}
\textbf{Modelo} & \textbf{Familia} & \textbf{Prec} & \textbf{Rec} & \textbf{F1} & \textbf{Prec} & \textbf{Rec} & \textbf{F1} & \textbf{AUC} \\
\midrule
XGBoost (Base) & Clásico & 0.85 & 1.00 & 0.92 & 0.00 & 0.00 & 0.00 & 0.6120 \\
LightGBM & Clásico & 0.85 & 1.00 & 0.92 & 0.15 & 0.001 & 0.01 & 0.6250 \\
CatBoost 5to Intento & Clásico & 0.85 & 1.00 & 0.92 & 0.22 & 0.001 & 0.01 & \textbf{0.6315} \\
\midrule
Self-Supervised Masking & Locura & 0.85 & 1.00 & 0.92 & 0.00 & 0.00 & 0.00 & 0.6280 \\
UMAP + Clustering & Locura & 0.85 & 1.00 & 0.92 & 0.00 & 0.00 & 0.00 & 0.6303 \\
PyTorch NN + SMOTE & Locura & 0.86 & 0.89 & 0.88 & 0.25 & 0.196 & 0.22 & 0.6033 \\
CatBoost + Gaussian Noise & Locura & 0.85 & 1.00 & 0.92 & 0.00 & 0.00 & 0.00 & 0.6293 \\
\midrule
\rowcolor{green!10}
\textbf{Campeón (Tribu + Pesos)} & \textbf{Ensamblado} & \textbf{0.88} & \textbf{0.64} & \textbf{0.74} & \textbf{0.20} & \textbf{0.531} & \textbf{0.29} & \textbf{0.6305} \\
\bottomrule
\end{tabular}
}
\caption{Comparativa detallada de métricas. Note que el "CatBoost 5to Intento" alcanzó el máximo AUC (0.6315) pero siendo completamente ciego a la Clase 1 (Recall 0.001). El Campeón sacrifica 0.001 de AUC a cambio de atrapar al 53.1\% de los compradores.}
\end{table}

\subsection{Análisis del Trade-Off (0 vs 1)}
Como se observa en la tabla, casi todos los algoritmos (incluyendo el CatBoost que llegó a 0.6315) actuaron de forma "cobarde". Para maximizar el AUC global, el modelo predecía casi siempre "0" (logrando Precision y Recall perfectos en la Clase 0), e ignorando por completo a la Clase 1 (Recall 0\%). 

Nuestro \textbf{Modelo Campeón} intervino quirúrgicamente esta asimetría. Al aplicar el "Kaggle Trick" (Pesos Balanceados), forzamos al modelo a fallar en los "0" (bajando su Recall de 1.00 a 0.64) para obligarlo a arriesgarse con los "1" (subiendo su Recall de 0.00 a 0.53). Todo esto sin dañar la métrica global de AUC ($0.6305$).

\section{Arquitectura del Modelo Campeón}
La solución final no radica en una red neuronal incomprensible, sino en un pipeline tabular avanzado. A continuación se presenta gráficamente el flujo de procesamiento capa por capa:

\vspace{1em}
\begin{center}
\begin{tikzpicture}[
    node distance=2cm and 0cm,
    layer/.style={rectangle, rounded corners, draw=black, thick, fill=blue!10, text width=12cm, align=center, minimum height=1.5cm},
    arrow/.style={thick,->,>=stealth}
]

% Nodes
\node (layer1) [layer, fill=gray!20] {\textbf{Capa 1: Ingesta de Datos Brutos}\\ \small Cliente aislado: $X = \{\text{ingresos}, \text{saldo}, \text{edad}, \text{region}, \text{ocupacion}\}$};

\node (layer2) [layer, fill=cyan!10, below=of layer1] {\textbf{Capa 2: Ingeniería de Contexto Social (Tribu)}\\ \small Agrupamiento espacial por \texttt{region} + \texttt{ocupacion}.\\ Se calculan medias poblacionales ($\mu$) y se transforma al cliente en variables relativas: \\ $R_{ingreso} = \text{ingreso}_i / \mu_{ingreso}$ \\ $D_{saldo} = \text{saldo}_i - \mu_{saldo}$};

\node (layer3) [layer, fill=orange!20, below=of layer2] {\textbf{Capa 3: Motor de Árboles Oblivious (CatBoost)}\\ \small Construcción de árboles simétricos de profundidad 6. \\ Las variables categóricas son codificadas internamente usando Target Encoding ordenado.};

\node (layer4) [layer, fill=green!20, below=of layer3] {\textbf{Capa 4: Función de Pérdida Asimétrica (Kaggle Trick)}\\ \small LogLoss penalizada: $\mathcal{L} = - \sum \left( w_1 y \log(\hat{y}) + w_0 (1-y) \log(1-\hat{y}) \right)$ \\ Donde $w_1 = 85\%$ y $w_0 = 15\%$. Multa severa por ignorar la Clase 1.};

% Arrows
\draw [arrow] (layer1) -- (layer2);
\draw [arrow] (layer2) -- (layer3);
\draw [arrow] (layer3) -- (layer4);

\end{tikzpicture}
\end{center}
\vspace{1em}

\subsection{Desglose de Capas}
\begin{itemize}
    \item \textbf{Capa 1 (Datos Brutos):} Análisis determinó que las variables aisladas carecían de correlación lineal con el objetivo ($\rho < 0.07$).
    \item \textbf{Capa 2 (Contexto Social):} Resuelve el solapamiento. Al comparar matemáticamente al individuo con su "tribu", el modelo detecta clientes que "destacan" en su estrato socioeconómico local, revelando la intención de compra.
    \item \textbf{Capa 3 (CatBoost):} Utilizado por su manejo nativo de datos categóricos y resistencia al sobreajuste mediante esquemas de permutación.
    \item \textbf{Capa 4 (Pesos Asimétricos):} La intervención matemática final que curó la "ceguera" del modelo, multiplicando el error asociado a la clase minoritaria para equilibrar la actualización de gradientes.
\end{itemize}

\section{Conclusión Definitiva}
Las pruebas rigurosas demostraron que las técnicas de generación sintética (SMOTE, Ruido Gaussiano) fracasan en espacios financieros densamente solapados, creando "clientes imposibles" y degradando el AUC a $\sim0.60$. 

El óptimo global de este dataset reside en la \textbf{Ingeniería de Características Sociales sumado a una Función de Pérdida Asimétrica}. El modelo entregado captura la complejidad no lineal del comportamiento bancario, maximizando el valor de negocio mediante un asombroso 53\% de Recall en la detección de oportunidades reales.

\end{document}
